\Needspace{8\baselineskip}\originalchapter{2002 年入门作业：参数、屏蔽和导引中心}
本编处理 22.611J Fall 2002 原始作业中与已选题不完全相同的题。每题同时给原编号、条件和独立中文推导；原课程的手写答案另以“有官方答案”注明，不能把下文独立推导视为扫描原文。扫描中使用 Gaussian 单位时，下文显式转换成 SI。PS1 原题1、5 完整重复，交叉索引见本编末章。

\originalsection{2002 PS1：常数、Saha 与双温屏蔽}
\begin{exercise}\label{ex:hist02-1-2}
PS1 原题2（有官方答案）。查到两位有效数字并记忆：光速、电子电荷、电子与质子静质量能量（MeV）、1 eV 对应 K、空气与水分子数密度、由氢原子基态模型计算的电离能、Gauss/Tesla、每立方厘米／米的密度换算，以及 $\hbar c$ 与精细结构常数 $\alpha$。可任选 SI 或 Gaussian 单位并标明。
\end{exercise}
\begin{solution}
编者独立核算。取 SI：$c=3.0\times10^8\,\mathrm{m/s}$，$e=1.6\times10^{-19}\,\mathrm C$，$m_ec^2=0.51\,\mathrm{MeV}$，$m_pc^2=9.4\times10^2\,\mathrm{MeV}$，$1\,\mathrm{eV}/k_B=1.2\times10^4\,\mathrm K$。在约 $300\,$K、$1\,$atm 的空气，$n=p/(k_BT)\simeq2.4\times10^{25}\,\mathrm{m^{-3}}$；液水 $n=\rho N_A/M\simeq3.3\times10^{28}\,\mathrm{m^{-3}}$。Bohr 模型 $E_I=m_ee^4/[2(4\pi\epsilon_0)^2\hbar^2]=13.6\,$eV。$1\,\mathrm T=10^4\,\mathrm G$，$1\,\mathrm{cm^{-3}}=10^6\,\mathrm{m^{-3}}$，$\hbar c=197\,\mathrm{eV,nm}=1.97\times10^{-5}\,\mathrm{eV,cm}$，$\alpha=e^2/(4\pi\epsilon_0\hbar c)\simeq1/137$。空气值依题目未指定的温压而变，不能把上述参考条件视为唯一值。
\end{solution}

\begin{exercise}\label{ex:hist02-1-3}
PS1 原题3（有官方答案）。总压为一标准大气压的氢放电气体若服从题给 Saha 电离平衡，氢电离能为 $13.6\,$eV，Debye 长度远小于放电尺寸；估算一半氢核已电离时的温度。
\end{exercise}
\begin{solution}
编者独立解答。设氢核总密度 $N=n_H+n_p$，电离份额 $x=n_p/N=n_e/N=1/2$。理想混合气体 $p=(n_H+n_p+n_e)k_BT=(1+x)Nk_BT$。令题给 Saha 右端为 $S(T)=g(2\pi m_ek_BT/h^2)^{3/2}\exp[-E_I/(k_BT)]$，$g$代表采用的配分函数统计权重；左端 $n_en_p/n_H=x^2N/(1-x)=p/(3k_BT)$。因此解
\begin{equation}g\left(\frac{2\pi m_ek_BT}{h^2}\right)^{3/2}e^{-13.6\mathrm{eV}/k_BT}=\frac{p}{3k_BT}.
\end{equation}
采用常见氢基态权重 $g=1$，数值得 $T\simeq1.54\times10^4\,$K，$k_BT\simeq1.32\,$eV；若把电子自旋因子独立计为 $g=2$，则 $T\simeq1.46\times10^4\,$K。结果远低于电离能对应温度，是平动状态密度补偿 Boltzmann 因子的结果。需满足局域热平衡；仅有 $\lambda_D\ll$ 尺寸不足以保证 Saha 适用。
\end{solution}

\begin{exercise}\label{ex:hist02-1-4}
PS1 原题4（有官方答案）。电子与单价离子分别以 $T_e,T_i$ 服从 Boltzmann 因子，把点电荷 $Q$ 放入均匀等离子体；在 $|e\phi|\ll T_e,T_i$ 近似下求径向势 $\phi(r)$ 及屏蔽尺度。冷离子极限是否等同“不动离子”？这里温度用能量单位。
\end{exercise}
\begin{solution}
编者独立解答。$n_e=n_0\exp(e\phi/T_e)$、$n_i=n_0\exp(-e\phi/T_i)$，线性化 Poisson 式得
\begin{equation}(\nabla^2-\lambda^{-2})\phi=-\frac{Q}{\epsilon_0}\delta^{(3)}(\mathbf r),\qquad \lambda^{-2}=\frac{n_0e^2}{\epsilon_0}\left(\frac1{T_e}+\frac1{T_i}\right).
\end{equation}
球对称 Green 函数因而给 $\phi=Qe^{-r/\lambda}/(4\pi\epsilon_0r)$。$T_i\ll T_e$ 时离子静态响应 $|\delta n_i|/n_0\sim |e\phi|/T_i$ 更强、$\lambda$ 更短；不动离子是令 $\delta n_i=0$ 的动力学假设，得到电子 Debye 长度，绝非简单取 $T_i\to0$。点电荷近旁小势线性条件必失效，上式只在远区使用。
\end{solution}

\originalsection{2002 PS2：拓扑、极化和镜间加速}
\begin{exercise}\label{ex:hist02-2-1}
PS2 原题1（有官方答案）。C 形永磁体的磁通穿过铝环形体，降至 $1\,$K 成超导态后把环远远向右拉开。图示两个似乎矛盾的解释：场线被环无限拉长；或环中感生电流维持原磁通、远处仅呈小偶极，磁体附近场几乎恢复原样。哪幅场线图合理？另一幅的拓扑论证哪里失效？
\end{exercise}
\begin{solution}
编者独立解答。超导环孔内的总穿链磁通近似守恒，外磁体与环的互感随距离降至零，于是持久电流趋向有限 $I\simeq\Phi_0/L$；环远场是偶极型，不可能在远距离仍以无穷大电流维持连到磁体的强场束。因此第二种远场图合理。$\nabla\cdot\mathbf B=0$禁止磁单极，不规定每条场线的全球身份；在远离环的区域可有 $\mathbf B=0$ 的瞬时零点／分界线，场线作为瞬时积分曲线可以改变连接关系。理想导体的冻结定理只约束导体物质内的磁通（或处处有相应理想流体的区域），真空里的线不能被当作附着在环上的物质细绳。具体动态拓扑还受移动边界及位移电流影响。
\end{solution}

\begin{exercise}\label{ex:hist02-2-2}
PS2 原题2（有官方答案）。均匀磁场中慢变横向电场 $\mathbf E_\perp(t)$ 使 $\mathbf E\times\mathbf B$ 漂移动能改变，然而该漂移垂直 $\mathbf E$。用功率守恒求附加极化漂移，并与题卷的 Gaussian 表式 $\mathbf v_{\rm pol}=c\dot{\mathbf E}_\perp/(\Omega B)$ 对照。
\end{exercise}
\begin{solution}
编者独立解答。取 $\Omega=qB/m$ 为带符号频率，SI 下 $\mathbf v_E=\mathbf E\times\mathbf B/B^2$，$\dot K_E=m\mathbf v_E\cdot\dot{\mathbf v}_E=m\mathbf E\cdot\dot{\mathbf E}/B^2$。要求 $q\mathbf E\cdot\mathbf v_{\rm pol}=\dot K_E$，即
\begin{equation}\mathbf v_{\rm pol}=\frac{m}{qB^2}\frac{\dd\mathbf E_\perp}{\dd t}=\frac1{\Omega B}\frac{\dd\mathbf E_\perp}{\dd t}\quad\text{(SI)}.
\end{equation}
Gaussian 下 $\mathbf v_E=c\mathbf E\times\mathbf B/B^2$，得到题给多一个 $c$ 的版本。此解只取平行 $\dot{\mathbf E}$ 的一阶部分；完全导引中心推导还须 $|\partial_t|\ll|\Omega|$。
\end{solution}

\begin{exercise}\label{ex:hist02-2-3}
PS2 原题3（有官方答案）。质子在相向移动的平面磁镜间俘获，镜比 $R=5$，初总能 $1\,$keV、中点 $v_\perp=v_\parallel$，两镜距 $L$、各以 $v_M=10\,$km/s 向中点移动。磁矩不变，弹性镜反射，每次反射所用时间忽略且 $L$ 可近似常量；估计逃逸前能量、反弹次数和时间，时间容许约两倍误差。
\end{exercise}
\begin{solution}
编者独立解答。初 $W_{\perp0}=W_{\parallel0}=0.5\,$keV，磁矩不变使中点 $W_\perp$ 不变。镜损失锥边界为 $W_\perp/(W_\perp+W_\parallel)=1/R$，故临界 $W_\parallel=4W_\perp=2\,$keV，总能 $2.5\,$keV。一次撞向迎面移动的镜后 $v_\parallel$ 增约 $2v_M$，因此
\begin{equation}N\simeq\frac{\sqrt{2(2\mathrm{keV})/m_p}-\sqrt{2(0.5\mathrm{keV})/m_p}}{2v_M}\simeq15.5.
\end{equation}
相邻镜碰撞之间行程约 $L$；连续近似 $t\simeq\int\dd v\,L/(2v_Mv)=L\ln(v_f/v_0)/(2v_M)=L\ln2/(2v_M)$。$L$ 用米时 $t\simeq3.47\times10^{-5}(L/\mathrm m)\,$s；起止相位和双镜反射约定产生题目允许的倍数差。镜的非绝热运动或 $L$ 显著缩短会破坏估算。
\end{solution}

\begin{exercise}\label{ex:hist02-2-4}
PS2 原题4（有官方答案）。Alcator C-Mod 外中平面参数 $R=0.66\,$m、$a=0.21\,$m、$B\simeq4.5\,$T、$n=4\times10^{14}\,\mathrm{cm^{-3}}$、$T_i=2\,$keV、径向 $E_r=1\,$kV/cm。按低 $\beta$ 估计离子的 $E\times B$、梯度 $B$、曲率漂移大小及其与热速之比；须说明质量／热速约定。
\end{exercise}
\begin{solution}
编者独立估算。取单价氢离子，$v_{thi}=\sqrt{2T_i/m_p}\simeq6.19\times10^5\,$m/s，$E_r=10^5\,$V/m。环向场在外侧 $|\nabla B|/B\simeq1/R$，按 $\langle W_\perp\rangle=T_i$、$\langle W_\parallel\rangle=T_i/2$ 有
\begin{align}
v_E&=E_r/B=2.22\times10^4\,\mathrm{m/s},&v_E/v_{thi}&=3.59\times10^{-2},\\
v_{\nabla B}&=T_i/(eBR)\simeq6.74\times10^2\,\mathrm{m/s},&v_{\nabla B}/v_{thi}&\simeq1.09\times10^{-3},\\
v_c&=2\langle W_\parallel\rangle/(eBR)\simeq6.74\times10^2\,\mathrm{m/s},&v_c/v_{thi}&\simeq1.09\times10^{-3}.
\end{align}
若用单个热速度粒子 $v_\parallel^2=v_{thi}^2$ 而非 Maxwell 平均，则曲率项为 $1.35\times10^3\,$m/s；须避免把两种约定的数字混用。给定密度用于检查 $\beta\sim2\mu_0n(T_e+T_i)/B^2$，若 $T_e\sim T_i$，约 $3.2\times10^{-2}$，与低 $\beta$ 近似相容。$\mathbf E$ 径向，漂移主要向与 $\mathbf B$ 垂直的切向。
\end{solution}
